\newpage
\thispagestyle{empty}
\chapter*{\sffamily Resumen}
\addcontentsline{toc}{chapter}{Resumen}%
%\begin{center}
%\textbf{\large \thesisname}
%\end{center}
% ------------------------------------------------------------------
% GUÍA PARA ESCRIBIR EL RESUMEN (borrar al terminar; máx. ~300 palabras,
% un solo párrafo, tiempo pasado/presente según UNAL, sin citas):
%   1. Problema y contexto (1-2 frases).
%   2. Objetivo general del trabajo (1 frase).
%   3. Metodología / diseño usado (1-2 frases): tipo de estudio,
%      participantes, instrumentos, papel de la IA generativa.
%   4. Resultados principales esperados u obtenidos (1-2 frases).
%   5. Conclusión o aporte del trabajo (1 frase).
% Traducir luego al inglés para el Abstract, sin agregar información nueva.
% ------------------------------------------------------------------
\par En la sede Palmira de la Universidad Nacional de Colombia, Cálculo
Diferencial y Matemáticas Básicas registran pérdidas altas, y los estudiantes
remitidos a tutoría personalizada por sus dificultades en ellas tienden a
resolver imitando procedimientos y dependen del modelado del tutor. La
inteligencia artificial generativa, ya presente en su forma de estudiar, puede
sostener ese razonamiento o reemplazarlo. Esta propuesta tiene como objetivo
diseñar, implementar y analizar una secuencia de tareas mediada por inteligencia
artificial generativa, orientada a la resolución de problemas de cálculo
diferencial en estudiantes de Cálculo Diferencial y Matemáticas Básicas remitidos
a la tutoría personalizada del programa de Grupos de Estudio Autónomo (GEA). El
estudio es cualitativo y adopta la lógica del experimento de enseñanza. La
secuencia se centra en el límite y organiza cada tarea en tres momentos: trabajo
sin IA, trabajo con Gemini en modo de aprendizaje guiado y validación, en la que
el estudiante compara lo que hizo con y sin la herramienta. Los datos provienen
de la observación de las sesiones, las hojas de trabajo, las conversaciones con
Gemini, los registros de la tutoría y una entrevista semiestructurada basada en
tareas con una muestra de cuatro a seis estudiantes. El análisis usa categorías
derivadas de la distinción entre razonamiento imitativo y creativo, de la
perspectiva instrumental y de la idea de inteligencia híbrida, con triangulación
de fuentes y de investigador. Se esperan tres productos: la secuencia de tareas
con su análisis a priori, un protocolo y un manual de uso de Gemini, y
narrativas analíticas sobre las estrategias, los errores y el tipo de apoyo que
los estudiantes encuentran en la IA. El aporte del trabajo está en mostrar bajo
qué condiciones de diseño la IA generativa acompaña el razonamiento de
estudiantes con dificultades en las matemáticas del primer año, en lugar de
sustituirlo.
\\[1cm]
\textbf{Palabras clave:} \palabrasclave.

%\newpage
%\thispagestyle{empty}
%\chapter*{\sffamily Abstract}
%\addcontentsline{toc}{chapter}{Abstract}%
%\begin{center}
%\textbf{\large \thesisnameeng}
%\end{center}
%\par Abstract text.
%\\[2cm]
%\textbf{Keywords:} \keywords

%Usar si es requerido, de no necesitarse favor comentar con % las lineas
\newpage 
\thispagestyle{empty}
%\chapter*{\sffamily Zusammenfassung}
%\addcontentsline{toc}{chapter}{Zusammenfassung}%
%\begin{center}
%\textbf{\large \thesisnamelang}
%\end{center}
%\par Zusammenfassung Texte.
%\\[2cm]
%\textbf{Schlüsselwörter:} \schlusselworter